\documentclass[11pt]{article}
\usepackage[a4paper,margin=28mm]{geometry}
\usepackage{amsmath,amssymb,amsthm,hyperref}
\newtheorem{theorem}{Theorem}
\newtheorem{corollary}[theorem]{Corollary}
\newtheorem{remark}[theorem]{Remark}
\title{Exact strength of the subset-divisibility obstruction}
\author{Research notes, 30 September 2026}
\date{}
\begin{document}
\maketitle
\noindent\textbf{Status.} The finite statements below have complete elementary proofs.
They concern necessary arithmetic conditions for permutation-fair dice, not
the existence of dice with prescribed face counts. No claimed resolution
of either principal size question follows.

Write $\vartheta(x)=\sum_{p\le x}\log p$, with sums over primes, and put
\[
 P_n=\prod_{p\le n}p^{n-p+1}
     =\exp\left(\sum_{j=1}^n\vartheta(j)\right).
\]
Call a positive integer vector $c=(c_1,\ldots,c_n)$ arithmetically admissible
if
\[
 |S|!\mid\prod_{i\in S}c_i\quad\hbox{for every }S\subseteq[n].
\]
Every permutation-fair face-count vector is arithmetically admissible.

\begin{theorem}[Complete characterization]
The following conditions on $c$ are equivalent:
\begin{enumerate}
\item $c$ is arithmetically admissible;
\item for every prime $p\le n$, at most $p-1$ coordinates of $c$ are not
divisible by $p$.
\end{enumerate}
Thus restrictions to prime-sized subsets already imply all the
subset-factorial divisibility conditions. Higher prime powers supply no
additional obstruction of this form.
\end{theorem}
\begin{proof}
Necessity follows by taking a set of $p$ coordinates not divisible by $p$.
For sufficiency, fix $S$ of size $k$ and a prime $p\le k$. At least
$k-p+1$ coordinates in $S$ are divisible by $p$. Legendre's formula gives
\[
 v_p(k!)=\frac{k-s_p(k)}{p-1}\le\frac{k-1}{p-1}\le k-p+1,
\]
where $s_p(k)\ge1$ is the base-$p$ digit sum; the last inequality is
equivalent to $(p-2)(k-p)\ge0$. Hence
$v_p(\prod_{i\in S}c_i)\ge v_p(k!)$. This holds for every relevant prime.
\end{proof}

\begin{theorem}[Sharp product obstruction, nearly sharp length obstruction]
Among arithmetically admissible vectors, the exact minimum of
$\prod_{i=1}^n c_i$ is $P_n$. If
\[
 B_n=\min\left\{\sum_i c_i:c\text{ arithmetically admissible}\right\},
\]
then
\[
 nP_n^{1/n}\le B_n\le n^2P_n^{1/n}.
\]
Moreover, an admissible squarefree vector attaining product $P_n$ can be
constructed greedily with $\max_i c_i/\min_i c_i\le n$.
\end{theorem}
\begin{proof}
The characterization forces each prime $p\le n$ into at least $n-p+1$
coordinates, so every product is divisible by $P_n$.
Start all coordinates at $1$. For each prime $p\le n$, multiply the
$n-p+1$ currently smallest coordinates by $p$. Each prime occurs in
exactly the required number of coordinates and nowhere with exponent
greater than $1$, so the final product is $P_n$ and the vector is
admissible.

If the old coordinates have maximum/minimum ratio at most $R$,
the new ratio is at most $\max(R,p)$. Indeed, for two coordinates that
are both multiplied or both left alone, their ratio is unchanged. The
ratio of a multiplied coordinate to an unmultiplied coordinate is at
most $p$, since all multiplied coordinates were among the smallest.
The reverse ratio is at most $R/p\le R$.
Induction proves the asserted ratio bound. Therefore every coordinate is
at most $nP_n^{1/n}$, giving the upper bound for $B_n$; the lower bound
is AM--GM.
\end{proof}

\begin{corollary}[Asymptotic barrier]
The prime number theorem gives
\[
 \log P_n=\left(\tfrac12+o(1)\right)n^2,
 \qquad
 B_n=\exp\left(\left(\tfrac12+o(1)\right)n\right).
\]
Consequently every permutation-fair collection has
\[
 \prod_i c_i\ge P_n,
 \quad\max_i c_i\ge P_n^{1/n},
 \quad\sum_i c_i\ge nP_n^{1/n}.
\]
No argument using only subset-factorial divisibility can improve the
coefficient $1/2$ in the exponential growth of the maximum or total
number of faces.
\end{corollary}
\begin{proof}
Use $\log P_n=\sum_{j=1}^n\vartheta(j)$ and
$\vartheta(j)=j+o(j)$. The explicit balanced admissible vectors prove
the barrier statement.
\end{proof}

For equal coordinates, arithmetic admissibility is equivalent to
$\prod_{p\le n}p\mid c_1$. Hence the exact equal-coordinate arithmetic
optimum is $n\exp(\vartheta(n))=\exp((1+o(1))n)$.
The arithmetic relaxation alone therefore exhibits a ratio
$\exp((1/2+o(1))n)$ between its equal and arbitrary optima. This is
compatible with a polynomial relation between actual dice optima,
but cannot establish one.

\begin{corollary}[Restriction quantiles]
If the face counts of a permutation-fair collection are arranged as
$c_{(1)}\le\cdots\le c_{(n)}$, then, for every $1\le r\le n$,
\[
 \prod_{i=1}^r c_{(i)}\ge P_r,
 \qquad c_{(r)}\ge P_r^{1/r}.
\]
In particular, as $r\to\infty$, $c_{(r)}\ge\exp((1/2+o(1))r)$.
\end{corollary}

\begin{remark}[Relation to existing literature]
Purcell's Proposition 1 records the equal-sided primorial obstruction:
\url{https://doi.org/10.1080/00029890.2024.2345576}.
The manuscript already proves that at least $n-p+1$ face counts must be
divisible by $p$. The new content here is the converse for the entire
arithmetic relaxation, its exact product optimum, and the matching
balanced-vector construction. Novelty beyond these checked sources has
not been established by an exhaustive literature search.
\end{remark}

\end{document}
